% 08-limitations.tex 中文对照版（译自 sections/08-limitations.tex；不加 \label）
% 四个冻结方向只在本节出现（论文规则）；framing-only 术语只作为“未实现”出现。
\section{局限}

\paragraph{证据尺度，一次性说明。} 五条限制界定本文的每一项主张，我们在这里一次说清，而不是让它们散落在各节。
(i)~\emph{参数}：所有参数的 provenance 都是 \texttt{guess}；\texttt{measured} 参数计数为零。
相变阈值（$0.70$/$0.85$）、趋势阈值（$0.02$）、积分增益、ACT-R 常数与权重阈值都是设计默认值，
因此没有任何阈值可以被称为已调优、已验证或最优。
(ii)~\emph{真实流量}：真实轮次为八，其中五轮有 ratio 观测，且 ratio 全为 $\approx 0$；
标定协议的准入条件——三十个有 ratio 观测的轮次、十个边界样本、二十个人工结果标签——未满足。
(iii)~\emph{小样本}：外部基线试点为每格五轮（单一宿主、单一模型路由、单日）；宿主 pilot 为十格、每格一次运行；
机制补充为每类十二个探针；补丁案例研究为二十条载荷；每套留出语料二十条。
(iv)~\emph{效应量}：除 ACT-R 类上一次预注册的符号检验（$12/12$，$p=2.4\times10^{-4}$）之外，
不报告置信区间或 $p$ 值，也没有测量遗忘率、幻觉率或端到端拦截率；凡用代理指标处均按代理指标标注。
本文的主张是可达性（reachability），不是效应量（effect size）。
(v)~\emph{泛化}：单一宿主、单一模型路由、token 由字符估算；这些都不能外推到另一个宿主、语言或 tokenizer。

\paragraph{外部基线与补丁回路的范围。} 三受试者比较（见“外部基线”一节）继承上文的小样本限制，
并且各受试者看到的输入形态不同：无插件宿主收到原始四轮对话记录，而记忆系统收到检索到的记忆。
Mem0 臂开箱即用的 \texttt{add} 路径是累加式的，因此除非显式删除，被取代的值会保留在存储中；
比较环境的限制见下文单独一段。
我们为 LINGSHU 报告的拦截是\emph{结构性}准入判定：闸门阻止注入并记录它，
但宿主仍会把用户消息送给模型，因此模型可能读到攻击文本。
这次试点还暴露了四个限定其读数含义的产品发现：事实提取器识别不了自然的纠正措辞（F-A），
并把最后的问题当成事实存下来（F-B）；工具网关把 schema 错误计为阻断，并在后续轮次注入提示（F-C）；
第四项，一轮注入滞后（F-D），是时序缺陷而不是设计极限：宿主在插件的 pre-step 状态更新之前就完成了拼装，
因此一轮的首次请求携带的是上一轮快照。它被定位到宿主生命周期，并在 v5.4 中通过以下方式修复：
在 \texttt{agent/inbox/claimed} 上预计算本回合状态、并重写首个被求值的段落；修复由 24 条独立回归断言钉住
（\texttt{lag-regressions.mjs}：真实钩子时序下 \texttt{anchor\_lag=false}、轮次幂等，以及保留旧路径作为回退）。
因此 S1 的模型层单元格改用修复后的重跑数据（见“外部基线”一节）。
红队案例研究（见“案例研究”一节）按构造就是样本内的——二十条载荷就是补丁集——
而四十条留出改写仍为 0/40，因为补丁在插件侧、留出评测在引擎侧。
补丁是实验室 fork，不是产品变更；它在良性流量上的误报代价未被测量；
产品架构（包括不存在任何自适应模块）未变。

\paragraph{沉浸式安全验证：多场景，样本仍小。} 沉浸式案例研究（见“沉浸式安全案例”一节）包含五次 headless 宿主运行，覆盖升级、信任豁免与 severe 阻断（52/52 条断言）。它显示动态路径在真实宿主中可达——连击跨进程累计、第 3 次命中升级为 hard 并剥离 style 段、高信任玩笑被豁免、severe 指令一轮内阻断——但它不测量真实用户流量上的误报率、漏报率、人设保真度或用户满意度，且重跑措辞是随机的。各档位背后的规则词表仍是桩，有八个开放盲区（见附录）。我们报告这项验证，是为了补上机制与真实宿主进程之间的那一环，而不是把它当作效应量估计。

\paragraph{合成数据与 mock 断言。} 回放集完全合成，且允许合成占比为一，这使它不足以用于部署决策。
插件证据是 mock 宿主断言（$188 + 148 + 45 + 15 = 396$）加上上述十格 pilot，而不是持续的真实流量。
触发词表以两种语言为主，token 用量由字符数估算，而不是由 tokenizer 产生。

\paragraph{混沌注入的范围。} 混沌故障是合成的单次注入（每种故障 $n=1$），不是生产故障的样本。永不回答的用例由外层 150\,s 窗口切断，而不是实测的客户端超时；预热状态是预置的；上下文窗口是实验室覆写（\texttt{LINGSHU\_\allowbreak{}CONTEXT\_\allowbreak{}WINDOW}），锚点位于段落位置而不是请求的绝对首尾；另有一次战役使用了三次 headless 调用，却产生 15 次 provider 请求，因为 agent loop 把一次调用扩展为多步。原始转录与逐字段证据文件均已保留，第~\ref{sec:exp:chaos} 节中每条主张都绑定到字段，而不是一句总结。

\paragraph{外部基线环境限制。} 官方 \texttt{mem0ai} 2.2.1 库使用的是本地非神经哈希二元组嵌入器（当时无可用嵌入 API），其默认配置在跨语言查询上产生空召回。这些是我们比较环境的限制，不是对 Mem0 的主张。

\paragraph{消融的功效。} 预注册的三臂消融现已执行（零模型、确定性）。
其原始设计只能支持粗粒度结论：基线—系统之间的保持率差异是定义性的，
十二个探针只能检出约 0.5 或更大的差异。一个补充实验把 ACT-R 探针扩到十二个（见“机制级消融”一节），
此时效应变得显著（12/12，$p=2.4\times10^{-4}$），而有两个机制臂按构造无法与闸门分离。
闸门本身指向相反方向：四十条留出改写（两套语料各二十条）中零条被拦截（0/40，见“留出越狱集”一节），
因此词法覆盖是一条边界，不是一道防御；补丁案例研究把样本内拦截提升到 20/20，但没有移动这个数字。

\paragraph{实现债。} 若干机制已实现并有单元测试，但未接线：僵局驱动的子目标化没有调用方；
降维路径没有冲突检测器，因此其触发条件恒为 \texttt{none}；UCA-03 不可达（拒绝标志从未被设置）；
动态权重淘汰从未被实际执行（回放预算从未超限）；饱和到闸门的连线由调用方提供。
拼装出的 prompt 不含用户原文，因此用户轮次的越狱无法在引擎侧被捕获；ACT-R 只实现了 base-level 激活。

\paragraph{解释性框架不是实现。} 有两个学术透镜以 \emph{framing-only}（仅作框架）方式出现：
生存系统模型（Viable System Model）与率失真（rate-distortion）类比。
二者都没有代码落点、参数或依赖：实现中不存在 VSM 这一术语，也没有计算任何率失真泛函。
它们只用于定位本工作，不能作为证据引用。

\paragraph{冻结的研究方向。} 四个方向被明确\emph{冻结}（frozen）、\emph{未实现}（not implemented），
且\emph{没有任何工程证据}（no engineering evidence）：(i) 行动闭环（带副作用回滚的闭环工具调用）；
(ii) 自适应演化（重写自身搜索目标的元学习）；(iii) 群体心智（多 Agent 黑板）；
以及 (iv) 多模态感知。它们都没有原型、依赖或配置开关。
它们的缺席界定了我们的主张：控制面管辖 prompt 拼装，不管辖工具执行、跨 Agent 协调或非文本输入。

\paragraph{安全范围。} fail-closed（失败即阻断）保证是行为性的，且只覆盖注入侧。
我们没有实现独立的危机安全前置层；拒绝依赖宿主安全层与 prompt 级宪法。
因此我们不作任何端到端安全主张：“外部基线”一节的结构性比较测量的是准入判定（闸门裁决与审计事件）
加上一个小规模模型侧 pilot，而不是一个可迁移到其他宿主、模型或攻击分布的安全率。
